\documentclass[12pt,letterpaper]{article}
\usepackage[spanish]{babel}
\usepackage{amsmath}
\usepackage{amssymb}
\usepackage{hyperref}
\usepackage{listings}
\usepackage{xcolor}
\usepackage{geometry}
\usepackage{graphicx}
\usepackage{array}
\usepackage{booktabs}

% Configuración de código
\geometry{margin=1in}
\lstset{
    language=Python,
    basicstyle=\ttfamily\small,
    keywordstyle=\color{blue}\bfseries,
    commentstyle=\color{gray},
    stringstyle=\color{red},
    breaklines=true,
    numbers=left,
    numberstyle=\tiny,
    frame=single,
    backgroundcolor=\color{white},
}

\title{\textbf{Entrega 1: Esquema SRTBOT\\Problema de la Torta de Cumpleaños}}
\author{\textbf{Integrantes:} Benjamín Farías y Francisco Solís}
\date{\today}

\begin{document}

\maketitle

\section*{Resumen Ejecutivo}

En esta entrega se presenta el esquema \textbf{SRTBOT} (Subproblema, Relación, Tabla, Borde, Orden, Transición) como marco metodológico para resolver el problema de la torta de cumpleaños mediante programación dinámica. Se implementa utilizando tablas de hash (diccionarios en Python) con memoización top-down. El objetivo es establecer claramente la estrategia de solución antes de optimizaciones posteriores.

\section*{Alcance}

Este informe cubre:

\begin{itemize}
    \item \textbf{Análisis del problema}: Descripción formal del problema de la torta de cumpleaños y su modelado como un juego de suma cero con dos jugadores
    \item \textbf{Esquema SRTBOT}: Aplicación sistemática de los seis componentes (Subproblema, Relación, Tabla, Borde, Orden, Transición) al problema específico
    \item \textbf{Estrategia minimax}: Implementación de la lógica que permite al profesor maximizar su satisfacción mientras la hermana intenta minimizarla
    \item \textbf{Memoización con diccionarios}: Técnica de programación dinámica top-down usando tablas de hash para evitar recálculos
    \item \textbf{Análisis de complejidad}: Cálculo formal de las complejidades temporal ($O(n^3)$) y espacial ($O(n^2)$)
    \item \textbf{Pseudocódigo e implementación}: Presentación clara del algoritmo antes de su implementación en Python
\end{itemize}

El informe \textbf{NO} cubre:

\begin{itemize}
    \item Optimizaciones posteriores o versiones paralelas del algoritmo
    \item Comparación con otras técnicas de resolución (por ejemplo, programación dinámica bottom-up con arreglos multidimensionales)
    \item Problemas variantes del juego de la torta (por ejemplo, porciones con pesos diferentes, número impar de jugadores)
    \item Aplicaciones prácticas fuera del contexto académico
\end{itemize}

%===============================================================================
\section{Contextualización del Problema}
%===============================================================================

\subsection{Descripción del Problema}

El profesor y su hermana comparten una torta redonda dividida en $2n$ porciones iguales, numeradas de $0$ a $2n-1$ en sentido antihorario. Cada porción $i$ tiene asociado un valor de satisfacción $st_i \in \mathbb{Z}$ (puede ser negativo, cero o positivo).

\begin{itemize}
    \item Los dos jugadores se alternan para elegir un ángulo válido $\alpha_i$
    \item El jugador que elige $\alpha_i$ se come las porciones desde $\alpha_i$ hasta $\alpha_i + \pi$ (en sentido antihorario)
    \item Esto equivale a consumir exactamente $n$ porciones consecutivas
    \item El profesor comienza el juego (turno 0)
    \item Ambos jugadores juegan de manera óptima para maximizar su propia satisfacción
\end{itemize}

\subsection{Objetivo}

Determinar la máxima satisfacción total \textbf{garantizada} que puede obtener el profesor si ambos jugadores juegan óptimamente. En teoría de juegos, esta es la solución minimax del juego.

%===============================================================================
\section{Esquema SRTBOT}
%===============================================================================

\subsection{S: Subproblema}

Se define la función:
\[
f(i, j, t) = \text{máxima satisfacción garantizada para el Profesor}
\]

donde:
\begin{itemize}
    \item $i, j \in \{0, 1, \ldots, 2n-1\}$: índices inicial y final del rango de porciones dis:ponibles
    \item $j - i + 1 \geq n$: el rango debe contener al menos $n$ porciones
    \item $t \in \{0, 1\}$: indicador de turno
    \begin{itemize}
        \item $t = 0$: es turno del Profesor (maximiza su satisfacción)
        \item $t = 1$: es turno de la Hermana (minimiza la satisfacción del Profesor)
    \end{itemize}
\end{itemize}

\textbf{Interpretación:} El valor $f(i, j, t)$ representa la máxima satisfacción que el profesor puede garantizarse si ambos juegan óptimamente desde el estado actual (porciones $[i, j]$ disponibles, turno $t$).

\subsection{R: Relación de Recurrencia}

Sean $n$ = número de porciones por jugada, $\text{suma}(a, b)$ = $\sum_{k=a}^{b} st_k$ la suma de satisfacciones en el rango $[a, b]$.

\subsubsection{Casos Base}

\begin{align}
\text{Si } j - i + 1 &< n: &f(i, j, t) &= 0 \label{eq:base1}\\
\text{Si } j - i + 1 &= n: &f(i, j, t) &= \text{suma}(i, j) \label{eq:base2}
\end{align}

En el caso base \eqref{eq:base2}, cuando quedan exactamente $n$ porciones, solo es posible una última jugada que se las lleva todas. No hay turnos posteriores.

\subsubsection{Caso Recursivo}

Para $j - i + 1 > n$, consideramos todas las jugadas válidas del jugador actual:

\textbf{Si } $t = 0$ \textbf{(turno del Profesor — maximiza):}
\[
f(i, j, 0) = \max_{k=i}^{j-n+1} \left[ \text{suma}(k, k+n-1) + f(k+n, j, 1) \right]
\]

El profesor elige el índice $k$ donde comenzará su jugada, eligiendo $n$ porciones consecutivas desde $k$ hasta $k+n-1$, obteniendo $\text{suma}(k, k+n-1)$. Las porciones restantes $[k+n, j]$ quedan para el siguiente turno (hermana).

\textbf{Si } $t = 1$ \textbf{(turno de la Hermana — minimiza):}
\[
f(i, j, 1) = \min_{k=i}^{j-n+1} \left[ f(k+n, j, 0) \right]
\]

La hermana elige dónde comenzar su jugada de manera que minimice la satisfacción futura garantizada del profesor. La hermana no obtiene satisfacción directa en esta formulación; solo reduce las porciones disponibles.

\subsection{B: Borde (Casos Base)}

Los casos base determinan cuándo termina la recursión:

\begin{table}[h]
\centering
\begin{tabular}{|p{3cm}|p{2.5cm}|p{5cm}|}
\hline
\textbf{Condición} & \textbf{Valor} & \textbf{Justificación} \\
\hline
Porciones restantes $< n$ & $0$ & No es posible hacer una jugada válida \\
\hline
Porciones restantes $= n$ & $\text{suma}(i, j)$ & Última jugada, se llevan todas las porciones \\
\hline
\end{tabular}
\end{table}

\subsection{T: Tabla / Técnica de Memoización}

Se implementa utilizando una estructura de diccionario para memoización:

\subsubsection{Estructura de Datos: Tablas de Hash (Diccionarios)}

Estructura mediante diccionario:
\[
\text{memo}[(i, j, t)] = f(i, j, t)
\]

donde la clave es una tupla $(i, j, t)$ y el valor es la solución memoizada. Esta estructura permite almacenar solo los subproblemas que efectivamente se calculan, sin necesidad de preasignar espacio para todas las combinaciones posibles.

\subsection{O: Orden de Cálculo}

Se utilizará memoización con cálculo \textbf{bajo demanda} (top-down):

\begin{enumerate}
    \item Se inicia con la llamada a $f(0, 2n-1, 0)$ (profesor elige primero con todas las porciones)
    \item Cada llamada recursiva verifica si el resultado ya está memoizado
    \item Si no está, se calcula recursivamente
    \item El resultado se almacena en la tabla de memoización antes de retornar
    \item La recursión termina cuando se alcanza un caso base
\end{enumerate}

El conjunto de subproblemas visitados depende del árbol de recursión, pero en el peor caso se pueden visitar $O((2n)^2)$ pares $(i, j)$ con 2 posibles turnos.

\subsection{T (Segunda T): Transición y Complejidad}

\subsubsection{Transición de Estado}

Cada estado se calcula explorando todas las transiciones posibles:
\begin{itemize}
    \item Un estado $(i, j, t)$ transiciona a múltiples estados $(i', j', t')$
    \item Número de transiciones por estado: $O(n)$ (una por cada posición donde se puede hacer una jugada)
    \item Total de transiciones: $O((2n)^2 \cdot n) = O(n^3)$
\end{itemize}

\subsubsection{Complejidad Temporal y Espacial}

\begin{itemize}
    \item \textbf{Tiempo}: $O(n^3)$ en el caso promedio (para memoización)
    \begin{itemize}
        \item $O((2n)^2)$ subproblemas posibles
        \item $O(n)$ operaciones por subproblema
    \end{itemize}
    \item \textbf{Espacio}: $O(n^2)$ en promedio
    \begin{itemize}
        \item Profundidad de recursión: $O(n)$
        \item Tabla de memoización: $O((2n)^2) = O(n^2)$ entradas
    \end{itemize}
\end{itemize}

%===============================================================================
\section{Implementación de la Estrategia}
%===============================================================================

\subsection{Estructuras de Datos}

Se utiliza una estructura de memoización basada en diccionarios:

\subsubsection{Diccionarios (Hash Table)}

Estructura clara y eficiente para esta implementación:

\begin{lstlisting}
memo = {}  # Diccionario vacío
# Clave: (i, j, turno) - tupla con índices y turno
# Valor: f(i, j, turno) - máxima satisfacción garantizada
\end{lstlisting}

\subsection{Pseudocódigo del Algoritmo}

\begin{lstlisting}
function satisfaccion_hash(st, n):
    // Inicialización
    total_porciones ← 2 * n
    memo ← {}  // Diccionario para memoización

    // Función auxiliar: suma de rango [inicio, fin]
    function suma_rango(inicio, fin):
        if inicio ≤ fin:
            return sum(st[inicio : fin+1])
        else:
            return 0

    // Función recursiva principal
    function resolver(i, j, turno):
        // PASO 1: Verificar memoización
        clave ← (i, j, turno)
        if clave ∈ memo:
            return memo[clave]

        // PASO 2: Calcular tamaño del subproblema
        tamaño ← j - i + 1

        // CASO BASE 1: Menos de n porciones
        if tamaño < n:
            resultado ← 0

        // CASO BASE 2: Exactamente n porciones (última jugada)
        else if tamaño == n:
            resultado ← suma_rango(i, j)

        // CASO RECURSIVO: Más de n porciones
        else:
            if turno == 0:  // TURNO DEL PROFESOR (MAXIMIZA)
                mejor ← -∞

                for k from i to j-n+1:  // Posiciones válidas de inicio
                    ganancia_profesor ← suma_rango(k, k+n-1)

                    siguiente_i ← k + n
                    siguiente_j ← j

                    if siguiente_i ≤ siguiente_j:
                        futuro ← resolver(siguiente_i, siguiente_j, 1)
                    else:
                        futuro ← 0

                    valor_total ← ganancia_profesor + futuro
                    mejor ← max(mejor, valor_total)

                resultado ← mejor

            else:  // TURNO DE LA HERMANA (MINIMIZA)
                mejor ← +∞

                for k from i to j-n+1:  // Posiciones válidas de inicio
                    siguiente_i ← k + n
                    siguiente_j ← j

                    if siguiente_i ≤ siguiente_j:
                        futuro ← resolver(siguiente_i, siguiente_j, 0)
                    else:
                        futuro ← 0

                    mejor ← min(mejor, futuro)

                resultado ← mejor

        // PASO 3: Memoizar resultado
        memo[clave] ← resultado
        return resultado

    // INICIALIZACIÓN: todas las porciones, turno del profesor
    return resolver(0, total_porciones-1, 0)
\end{lstlisting}

\subsection{Ejemplo de Ejecución}

\subsubsection{Entrada Ejemplo}

\[
n = 2, \quad st = [3, 1, 2, -1]
\]

Porciones: 4 (2n=4). Cada jugada consume 2 porciones consecutivas.

\subsubsection{Árbol de Recursión (Parcial)}

\begin{verbatim}
resolver(0, 3, 0)  [Profesor elige entre posiciones 0 y 2]
├─ Opción 1: come [0,1] → suma=4
│  └─ resolver(2, 3, 1)  [Hermana elige posición 2]
│     └─ come [2,3] → suma=1
│        └─ Total Profesor: 4 + 1 = 5
├─ Opción 2: come [1,2] → suma=3
│  └─ resolver(3, 3, 1)  [Menos de 2 porciones]
│     └─ No puede jugar → 0
│        └─ Total Profesor: 3 + 0 = 3
└─ Opción 3: come [2,3] → suma=1
   └─ resolver(4, 3, 1)  [Rango inválido]
      └─ No puede jugar → 0
         └─ Total Profesor: 1 + 0 = 1

Máximo: max(5, 3, 1) = 5 (con opción 1)

Pero en opción 1, es turno de la Hermana (minimiza):
└─ resolver(2, 3, 1) [Hermana elige posición 2]
   └─ Opción única: come [2,3] → suma=1
      └─ resolver(4, 3, 0) → inválido, retorna 0
         └─ Total Profesor: 0
      Hermana minimiza: min(0) = 0

Entonces opción 1 da: 4 + 0 = 4

Resultado final: max(4, 3, 1) = 4
\end{verbatim}

\section{Ejemplos de Casos Base}

\begin{table}[h]
\centering
\begin{tabular}{|p{2cm}|p{3cm}|p{2cm}|p{2cm}|}
\hline
\textbf{$n$} & \textbf{st (satisfacciones)} & \textbf{Porciones} & \textbf{Resultado} \\
\hline
2 & [1, 2, 3, 4] & 4 & 6 \\
\hline
2 & [5, -5, 3, -3] & 4 & 2 \\
\hline
1 & [10, -5, 8, -2] & 4 & 10 \\
\hline
\end{tabular}
\end{table}

\section{Conclusiones}

El esquema SRTBOT proporciona una estructura clara y sistemática para abordar el problema de la torta de cumpleaños:

\begin{enumerate}
    \item La definición del subproblema captura el conflicto entre dos jugadores en un juego alternado
    \item La relación de recurrencia implementa la lógica minimax (maximizar para el profesor, minimizar para la hermana)
    \item La memoización con diccionarios permite evitar recalcular subproblemas de manera eficiente
    \item La complejidad resultante $O(n^3)$ es viable para valores moderados de $n$
    \item La implementación con tablas de hash es clara, intuitiva y fácil de depurar
\end{enumerate}

%===============================================================================
\section{Referencias}

\begin{itemize}
    \item \href{https://github.com/Zelechos/Introduccion_a_los_Algoritmos/blob/main/Introduction_to_Algorithms_Third_Edition_(2009).pdf}{Algoritmos de Programación Dinámica - Cormen, Leiserson, Rivest, Stein}
    \item Teoría de Juegos Combinatorios
    \item Técnica SRTBOT - Análisis de Algoritmos (Curso ADA)
\end{itemize}

\end{document}

